\documentclass[11pt]{article}
\usepackage[utf8]{inputenc}\usepackage[T1]{fontenc}\usepackage{lmodern}
\usepackage[margin=1in]{geometry}
\usepackage{microtype}\setlength{\emergencystretch}{3em}
\usepackage{amsmath,amssymb,amsthm,booktabs,graphicx,array,xcolor,enumitem,caption}

\usepackage[round]{natbib}
\usepackage[hidelinks]{hyperref}
% AUTO-GENERATED by code/make_macros.py. Do not edit by hand.
% Every value traces to the results file named in the comment.

   % measured values only

\newtheorem{definition}{Definition}
\newtheorem{proposition}{Proposition}

\title{\bf Coherent Errors: Chess Transformers Act\\on Misremembered Boards}

  \author{
    Mohit Pratap Singh Rathore\\ \small Oviqo Private Limited
    \and
    Gunveer Sinh Kalsi\\ \small Oviqo Private Limited
    \and
    Sriharsha Meduri\\ \small Oviqo Private Limited
  }

\date{}

\begin{document}
\maketitle

\begin{abstract}
Long-horizon failure is reported as one number, which confounds two failures with
opposite remedies: the model has lost track of the world, or it holds the world
and chooses badly. We separate them where an exact state oracle exists, chess
from move notation with no search at inference, in models of
3,440,064 to 39,884,288 parameters, reading linear probes
against the true position at every token.

Aggregate state fidelity, the quantity usually reported, is the weaker predictor
of an illegal next move within a fixed depth, at AUC 0.538 to
0.671 against 0.592 to 0.687 for the same
representation restricted to the squares the move touches. And the errors are
coherent. Reconstructing the position the model appears to believe,
0.355 of its illegal moves are legal there, against
0.050 for the same move under another position's belief and
0.007 for an arbitrary illegal move. Editing the probe direction
encoding an occupied square removes 0.934 of the probability mass
on moves from it, a norm-matched random edit 0.077, so the
representation is read and not merely present. Belief consistency rises
monotonically across six model sizes; the causal edit is not measured at every
size. A substantial identified subset of long-horizon errors is coherent with a
degraded state representation that the policy demonstrably reads.

We report failures too: a backwards layer profile, a degenerate pre-registered
metric, an auxiliary objective improving the measurement of state more than its
consequences, and a second domain needing seven designs to become measurable.
\end{abstract}

\noindent\textbf{Keywords:} state tracking, world models, probing, model
diagnostics, chess, long-horizon reasoning

\section{Introduction}

A model that reasons over many steps has to keep track of something. An
attention-only decoder has no explicit bounded mutable state in which to keep
it. Each token's representation is
rebuilt from the whole prefix through attention (Vaswani et~al. 2017), so
any quantity the task requires, a board position, the value of a variable, which
sub-goal remains open, exists only because the network re-derives it, at every
layer and every position, from the beginning. That re-derivation is dependable
over short spans and unreliable over long ones. Anyone who has watched a model
reason for forty steps knows the symptom: it stays fluent and quietly stops being
right about its own earlier commitments.

The usual remedies sit outside the model. Tree search grafts deliberation onto
the policy (Yao et~al. 2023). Process reward models supervise intermediate
steps with a learned verifier (Lightman et~al. 2024). Self-consistency samples
several chains and votes (Wang et~al. 2023). Each raises end-task
accuracy, and each leaves the original question untouched: what fails inside the
network as the chain lengthens? The question matters commercially as well as
scientifically, because the two candidate failures recommend opposite purchases.
If the model has forgotten the world, a verifier that rescores its outputs is
paying repeatedly to repair something the architecture gave away. If the model
remembers the world and reasons badly over it, a verifier is exactly right.

Answering the question requires reading an intermediate representation against a
known-correct intermediate state, and on open-ended text no such state exists.
The intermediate ``state'' of an argument in prose is not well defined, which is
why evaluation in this area is noisy and why claims about it are hard to falsify.
Our response is to move the measurement to a domain where the state is exact,
cheap, and external, establish what can be learned there, and be explicit about
what that does and does not license.

\paragraph{What we find.}
Working with search-free autoregressive models trained only on chess move
notation, we measure the internal state against a rules engine at every token and
compare it with what the model actually does. Three results carry the argument, and the rest of
the paper tests how far they reach.

\begin{enumerate}[nosep]
\item \textbf{The errors are coherent.} We reconstruct the board the model
appears to believe and ask whether its illegal move is legal there. It is, for
0.355 of illegal moves, with a bootstrap interval of
[0.346, 0.362]. The same move scored against
another position's believed board at equal depth is legal 0.050 of
the time, and an arbitrary illegal move scored against the model's own believed
board is legal 0.007 of the time. Lossy decoding caps what is
detectable at 0.654, the rate for moves that were genuinely legal.
\item \textbf{The representation is read, not merely present.} Editing the probe
direction that encodes an occupied square removes 0.934 of the
probability mass on moves from that square, while a norm-matched random edit
removes 0.077 and shifts the mass at chance. The edit
strength is calibrated rather than chosen.
\item \textbf{Where the world model is measured matters, increasingly with
depth.} Whole-board probe accuracy predicts an illegal next move with an AUC of
0.538 to 0.671. The same representation restricted to
the squares the move uses reaches 0.592 to 0.687. In
the opening the two are almost indistinguishable; the gap opens as reasoning
deepens, and in a joint fit the move-touched term carries a standardised
weight of 0.504 against 0.118.
\end{enumerate}

The first two are the ones we would defend hardest, and they belong together:
the errors are coherent with a decayed state, and that state is causally
load-bearing rather than an epiphenomenon. Failure here is not noise breaking
through as the context grows. A substantial share of it is action admissible under a
world that has drifted, which is a claim about state rather than about reasoning
for that share, and it suggests why fluency outlives correctness: the model acts
coherently with respect to a position it no longer holds.

\paragraph{Why chess.}
Chess presented purely as move notation, with no tree search at inference, gives
what prose does not. The true state after any prefix is computable by a rules
engine, so probing has a per-token label. Legality is decidable, so a
state-tracking failure has an unambiguous observable rather than a judgement
call. Move quality is separately scoreable by an engine, so the second failure
mode can be measured conditional on the first being absent. Depth is uniformly
scaled, unlike ``number of reasoning steps'' in prose. And the domain is
combinatorial, so surface heuristics do not substitute for computation. Chess is
our instrument, not our subject, and we are explicit in
Section~9 about the reach that buys.

\paragraph{Contributions.}
\begin{enumerate}[nosep]
\item \textbf{A substantial identified subset of long-horizon errors is
coherent.} Invalid actions are valid in the position the model appears to believe
far more often than matched controls allow, and the identified share rises across
the range of model sizes we train.
\item \textbf{The state representation is read, not merely present.} Inducing a
false belief by editing a single probe direction changes the action accordingly,
at a calibrated edit strength and against a norm-matched control.
\item \textbf{Where a world model is measured matters.} Whole-board probe
accuracy predicts whether the next action will be valid, and the same measurement
restricted to the state that action consumes predicts it better, with the gap
widening as reasoning deepens. We define a horizon on that fidelity directly,
which avoids a failure mode of behavioural horizons that we document.
\item \textbf{A per-token diagnostic stratification} into state mismatch,
belief-consistent error, conditional planning loss, and unresolved failure,
measured against an exact oracle rather than a learned verifier, with the
conditioning analysis that shows how far the strata can be told apart.
\item \textbf{An account of what failed}, including a degenerate metric, a
prediction that ran backwards, an intervention that improved the measurement more
than the behaviour, and seven task designs before a second domain became
measurable. Each changed a decision and none is usually reported.
\end{enumerate}

\section{Related Work}

\paragraph{Inference-time search.}
Deliberation over candidates (Yao et~al. 2023), process supervision
(Lightman et~al. 2024), and sampling with a vote
(Wang et~al. 2023), built on chain-of-thought prompting
(Wei et~al. 2022), all buy accuracy with inference compute and take the base
policy as given. Our concern is the base policy. A verifier that spends its
budget repairing state-tracking errors is compensating for a representational
defect at a recurring cost. What we do differently is diagnose which defect is
present before recommending a remedy.

\paragraph{Long-context evaluation.}
Position-dependent retrieval failure (Liu et~al. 2024) and effective-context
measurement (Hsieh et~al. 2024) show usable context is shorter than advertised
context. Both measure retrieval of material present in the prefix. Long-horizon
reasoning failure is different: the needed quantity was never written down, only
implied, and has to be maintained through composition. A model can attend to
every token of a sixty-ply game and still not know where the pieces are. What we
do differently is measure maintenance rather than retrieval.

\paragraph{World models and probing.}
Linear probes recovering board state from game transcripts
(Li et~al. 2023), the linear structure of those representations
(Nanda et~al. 2023), chess as a state-tracking testbed
(Toshniwal et~al. 2022), and causal editing of internal associations
(Meng et~al. 2022) together establish that sequence models form internal
world models and that intervening on them changes behaviour. Closest to our
setting, Karvonen (2024) probes board state and estimated player
skill in a chess-playing language model and intervenes on those representations,
and the same board-state probes are later used as ground truth for evaluating
dictionary learning (Karvonen et~al. 2024). That work establishes the
representation exists and is manipulable in this exact domain, which is our
premise rather than our result. That literature
answers whether a world model exists. We take its existence as given and ask a
different question: how the fidelity of that world model varies with reasoning
depth, and whether its failures explain the model's behavioural failures at the
level of individual actions. Probe design
follows the control-task discipline of Hewitt and Liang (2019).

\paragraph{Concurrent work, and what is actually new here.}
The measurement at the centre of this paper is not new.
Balogh and Jelasity (2026) report the fraction of illegal chess moves that
are legal on a probe-reconstructed board, and reach a conclusion close to the
opposite of ours: for models trained on large datasets that fraction is small, and
the gradient of their board-state probe is nearly orthogonal to that of the
next-token head, which they read as the probe lacking a causal role in
generation. Li and Jiang (2026) separate state tracking from decision quality
in searchless chess transformers, treating illegal moves as direct evidence of
tracking failure. Walker and Lyons (2026) benchmark exact state tracking across
architectures over a parameter range close to ours, and
Pereira and Zuidema (2026) show world-state representations decaying during
generation in frontier reasoning models, with a causal restoration experiment.
Harang et~al. (2025) document loss of coherent internal state over long
sequences. State decay with depth, chess state tracking, and scaling this
measurement are therefore not novel contributions of ours, and we do not claim
them.

What differs is the conditions and the conclusion. Their illegal moves are
adversarially induced; ours arise on ordinary held-out trajectories, which is a
different population and may have a different mechanism. We add same-depth
mismatched-belief and random-move controls, which bound how much of the effect a
merely plausible board would produce. We compare aggregate against
action-conditioned predictors on equal footing rather than assuming the aggregate
is the right one. And where their gradient-orthogonality analysis finds the probe
weakly coupled to generation, our calibrated residual-stream intervention finds a
strong effect on the action. Those two results are not strictly contradictory,
since a direction can be causally load-bearing at one layer while its gradient is
near-orthogonal to the output head, but the tension is real and we flag it as the
most useful open question this paper raises.

\paragraph{Expressivity of state.}
Theory bounds what these architectures can track: state-space models are weaker
than their recurrent form suggests (Merrill et~al. 2024), and intermediate
tokens change what a Transformer can express (Merrill and Sabharwal 2024). That
work asks what is representable in principle. We ask what a trained model
actually maintains, and where it stops.

\paragraph{Recurrence and state-space models.}
Structured state-space models (Gu et~al. 2022) and selective variants
(Gu and Dao 2024) give constant per-token state, usually motivated by
throughput and context length. We use a small supervised recurrent channel not as
a contribution in itself but as a test case for whether the diagnostic
discriminates between architectures.

\paragraph{Search-free chess models.}
Policies that reach strong play by distilling engine evaluations without a tree
(Ruoss et~al. 2024) show the task is learnable without search. Their
objective is strength; ours is measurement, so we deliberately study models weak
enough to fail. A model that never loses the position teaches nothing about
losing it.

\section{Design Problem and Requirements}

Stated plainly: long-horizon failure is reported as a single number, so it cannot
be acted on. A practitioner watching accuracy fall at step thirty cannot tell
whether to buy a larger model, add a verifier, or change the architecture,
because the measurement does not distinguish the causes.

The procedure rests on four commitments, each ruling out a way of fooling
oneself. The diagnostic is grounded in an external oracle rather than model
self-report. The two failure modes are separated before any effect is attributed.
Every decodability number is reported against trivial predictors that could
produce it without a world model. And the representation is tested by
intervention, not only by correlation.

\paragraph{In scope.} Search-free autoregressive decoders; domains with an exact,
cheaply computable state oracle; state of bounded description length.

\paragraph{Out of scope.} Inference-time search, which we neither use nor
evaluate; tasks whose state grows without bound; and claims about
production-scale models, which our compute does not license.

\section{Problem Formalisation}

Let a task be a token sequence $x_{1:L}$ generated left to right, and
$s_t = \Phi(x_{1:t})$ the exact state after the prefix, from an external oracle
$\Phi$. For chess $s_t = (b_t, \tau_t, c_t, e_t)$: piece placement, side to move,
castling rights, and the en-passant target.

Our probe decodes only the first component. Write
$\hat{b}_t^{(\ell)} = g^{(\ell)}(h_t^{(\ell)})$ for decoded piece placement, and
let the pseudo-state used by the belief test be
$\tilde{s}_t = (\hat{b}_t, \tau_t, \varnothing, \varnothing)$, with side to move
taken from ply parity and castling and en-passant disabled so that nothing is
licensed by an assumption we did not decode. Claims below concern
$\hat{b}_t$ and $\tilde{s}_t$, not $s_t$. The state space is large, but the description length of any
$s_t$ is constant in $t$: a position occupies a few hundred bits however many
plies produced it.

Let $h_t^{(\ell)}$ be the residual-stream activation at position $t$, layer
$\ell$, and $g^{(\ell)}$ a linear decoder trained to recover $s_t$ from it. Write
$\hat{s}_t^{(\ell)} = g^{(\ell)}(h_t^{(\ell)})$ for the decoded, or believed,
state.

\begin{definition}[Board-placement fidelity]
$F^{(\ell)}(t) = \Pr[\hat{b}_t^{(\ell)} = b_t]$, the probability that piece
placement is linearly decodable at depth $t$. We name it for placement rather
than for state because side to move, castling rights and the en-passant target
are not decoded.
\end{definition}

\begin{definition}[Move-touched fidelity]
The probability that $\hat{s}_t$ and $s_t$ agree on the origin and destination
squares of the selected move. We name it for what it measures rather than calling it action-relevant, because
it is a deliberately local proxy. The state an action
actually depends on is larger: legality can turn on intermediate path squares for
sliding pieces, the king's square and post-move attack status, pins and
discovered attacks, castling rights and transit squares, the en-passant target,
side to move, and promotion. A fuller treatment would define
$R_{\mathrm{legal}}(a, s)$, the complete set of components legality depends on,
and report a progression from whole-board through move-touched to path-aware and
fully legality-aware fidelity. We report the local proxy and the controls that
bound it, and leave that progression to future work.
\end{definition}

\begin{definition}[Behavioural error and its stratification]
Let $Y_t = 1$ when the action selected at depth $t$ is behaviourally invalid, and
let $M_t = 1$ when the decoded state disagrees with the oracle on the components
the selected action reads. By the law of total probability,
\[
\Pr(Y_t{=}1) = \Pr(M_t{=}1)\Pr(Y_t{=}1 \mid M_t{=}1)
             + \Pr(M_t{=}0)\Pr(Y_t{=}1 \mid M_t{=}0).
\]
This is an observational stratification and not a causal decomposition. The first
term is not by itself memory-caused error and the second is not by itself
policy-caused error.
\end{definition}

\begin{definition}[Belief-consistent error]
Write $\mathcal{A}(s)$ for the actions valid in state $s$ and
$\widetilde{\mathcal{A}}(\hat{s})$ for those admitted by the decoded state. An
error is \emph{belief-consistent} when
$B_t = \mathbf{1}[\,a_t \notin \mathcal{A}(s_t) \wedge a_t \in
\widetilde{\mathcal{A}}(\hat{s}_t)\,]$, and we report
$\Pr(B_t{=}1 \mid a_t \notin \mathcal{A}(s_t))$.
\end{definition}

Two consequences of this framing are worth stating plainly, because an earlier
version of this paper got them wrong. A state mismatch need not produce a
behavioural error: our own measurements show the probe wrong on the squares a
move touches for a substantial fraction of \emph{legal} moves, so
$M_t{=}1$ does not imply $Y_t{=}1$. And because the probe is an incomplete and
lossy read of the state, belief consistency identifies a subset of errors rather
than partitioning all of them. The errors it does not identify are
\emph{unresolved}: they may be genuine policy failures, state errors the probe
did not capture, information the probe cannot read linearly, or reconstruction
failures. We do not call them planning errors.

\section{The State-Space Bottleneck}

\begin{proposition}[Cost mismatch]
Under full attention each new token reads an $O(t)$ key-value history and
retains $O(t)$ cached memory, while the task state has description length
$O(1)$. A key-value cache does amortise the computation of earlier activations,
so the cost is not recomputation from scratch; what the architecture lacks is a
bounded mutable state that could carry $s_t$ forward and be updated in place. We
offer this as a motivating cost mismatch, not as a proof that fidelity must
decay.
\end{proposition}

Two consequences were testable and we tested both. Fidelity should fall with
depth even well inside the context window, because the re-derivation circuit must
compose more updates within a fixed layer budget. It does
(Section~7.1).

We also predicted fidelity would peak at intermediate layers and fall toward the
output, the state being a means rather than an end. \textbf{That prediction is
wrong.} Decodability rises through the stack and is maximal at the final block,
from 0.323 at the first layer to 0.543 at the
last. For this task the state sits close to the target: the next legal move is a
direct function of the position, so a representation making the position linearly
available is also a good representation for the output head. We keep the failed
prediction because it constrains where a state carrier should be inserted, and
our spacing was chosen before we knew.

\begin{figure}[t]
\centering
\includegraphics[width=6in]{layers.png}
\caption{State decodability by layer and depth. Contrary to our prediction, the
position is most linearly available at the final block rather than in the middle
of the stack.}

\end{figure}

\section{Diagnostic Protocol}

\paragraph{Data.} Public game records above a rating floor, deduplicated by
move-sequence hash, split by \emph{game}, so no position from a given game appears in more
than one split. Splitting by game does not make the splits prefix-disjoint,
since distinct games share openings, and we do not claim it does. The corpus holds 120,000 games:
96,000 for training, 12,000 for probes, 12,000
for evaluation, mean length 78.2 plies, move vocabulary
1,926.

\paragraph{Oracle.} A rules engine replays every game and emits the exact state
after every token, verified complete before any model is measured.

\paragraph{Probes.} One linear head per layer, fitted on the probe split and
reported on the evaluation split. Because most squares are empty, raw per-square
accuracy is misleadingly high, so we report accuracy on occupied squares against
a per-square, per-depth majority predictor.

\paragraph{Controls.} A majority predictor; a label-permutation control pairing
activations with another game's state at equal depth, preserving marginals while
destroying the relation to the input; and a randomised-weight model of identical
architecture, bounding what probe capacity alone extracts.

\paragraph{Behavioural measures.} Illegal top-1 move rate against depth, and the
probability mass on legal continuations, both read against the oracle.

\section{Results}

\subsection{State and behaviour both decay with depth}

\begin{table}[t]
\centering\small
\caption{State decodability against every trivial explanation, for the
baseline model. Occupied-square accuracy at the most decodable layer,
beside a per-square per-depth majority predictor, a label-permutation
control, and a randomised-weight model of identical architecture. The gap
is what training put there.}

\begin{tabular}{lccccc}
\toprule
Ply & model & majority & label perm. & random weights & illegal rate \\
\midrule
0--10 & 0.975 & 0.867 & 0.879 & 0.898 & 0.0004 \\
10--20 & 0.876 & 0.599 & 0.673 & 0.641 & 0.0115 \\
20--30 & 0.761 & 0.444 & 0.487 & 0.439 & 0.0439 \\
30--40 & 0.655 & 0.359 & 0.315 & 0.297 & 0.0982 \\
40--50 & 0.557 & 0.261 & 0.166 & 0.195 & 0.1766 \\
50--60 & 0.475 & 0.073 & 0.072 & 0.141 & 0.2347 \\
60--80 & 0.353 & 0.000 & 0.019 & 0.071 & 0.2920 \\
80--120 & 0.173 & 0.000 & 0.001 & 0.038 & 0.3231 \\
120--161 & 0.066 & 0.000 & 0.000 & 0.030 & 0.3192 \\
\bottomrule
\end{tabular}
\end{table}

\begin{table}[t]
\centering\small
\caption{Aggregate and move-touched state error compared like for
like, as predictors of the same binary outcome on the same samples:
whether the model's top-1 move is illegal. AUC is computed within each
depth bucket so neither predictor can simply be tracking depth. The
standardised coefficients come from a logistic fit containing both
predictors. The aggregate measure is informative and it is the weaker of
the two, and the gap widens with depth.}

\begin{tabular}{lccccr}
\toprule
Ply & AUC aggregate & AUC move-touched & $\beta$ aggregate & $\beta$ move-touched & $n$ \\
\midrule
10--20 & 0.671 & \textbf{0.687} & +0.065 & \textbf{+0.153} & 15{,}000 \\
20--30 & 0.627 & \textbf{0.681} & +0.120 & \textbf{+0.332} & 15{,}000 \\
30--40 & 0.584 & \textbf{0.662} & +0.100 & \textbf{+0.506} & 14{,}511 \\
40--50 & 0.548 & \textbf{0.643} & +0.057 & \textbf{+0.573} & 13{,}176 \\
50--60 & 0.551 & \textbf{0.617} & +0.114 & \textbf{+0.500} & 11{,}262 \\
60--80 & 0.538 & \textbf{0.592} & +0.076 & \textbf{+0.431} & 16{,}426 \\
\midrule
pooled, depth as covariate & n/a & n/a & +0.118 & \textbf{+0.504} & 85{,}375 \\
\bottomrule
\end{tabular}
\end{table}

\begin{table}[t]
\centering\small
\caption{Does state loss explain behaviour? The aggregate column is the
within-depth AUC of whole-board probe error as a predictor of an illegal
move, reported as a range over depth buckets. Move-touched error is the
probability the probe is wrong on the origin or destination square.
Belief consistency is the share of illegal
moves admissible under the model's own decoded board, beside the same
move under another position's belief, an arbitrary illegal move, and the
legal-move recovery reference.}

\begin{tabular}{lrccccccc}
\toprule
& & aggregate & \multicolumn{2}{c}{move-touched} & \multicolumn{4}{c}{belief consistency} \\
\cmidrule(lr){4-5}\cmidrule(lr){6-9}
Model & Params & AUC & illegal & legal & own & mismatch & random & recovery \\
\midrule
Attention-only & 7{,}345{,}920 & 0.54 to 0.67 & 0.806 & 0.384 & \textbf{0.355} & 0.050 & 0.007 & 0.654 \\
Attention-only, parameter-matched & 7{,}505{,}976 & n/a & 0.804 & 0.391 & \textbf{0.350} & 0.051 & 0.007 & 0.650 \\
Recurrent channel ($\lambda{=}0$) & 7{,}508{,}416 & n/a & 0.801 & 0.384 & \textbf{0.358} & 0.053 & 0.007 & 0.658 \\
Recurrent channel ($\lambda{=}1$) & 7{,}508{,}416 & n/a & 0.806 & 0.399 & \textbf{0.366} & 0.093 & 0.021 & 0.656 \\
6L/192 & 3{,}440{,}064 & 0.52 to 0.70 & 0.831 & 0.447 & \textbf{0.266} & 0.049 & 0.006 & 0.581 \\
12L/384 & 22{,}835{,}328 & 0.53 to 0.71 & 0.747 & 0.327 & \textbf{0.412} & 0.056 & 0.009 & 0.713 \\
12L/512 & 39{,}884{,}288 & 0.53 to 0.74 & 0.711 & 0.297 & \textbf{0.439} & 0.059 & 0.009 & 0.745 \\
\bottomrule
\end{tabular}
\end{table}

\begin{table}[t]
\centering\small
\caption{Causal intervention at the calibrated edit strength. Editing the
state direction removes probability mass from moves out of the affected
square; a norm-matched random direction does not. The manipulation check
reports how often the edit actually changed the decoded belief.}

\begin{tabular}{lccccc}
\toprule
Model & $n$ & belief flipped & mass before & mass after & random edit after \\
\midrule
Attention-only & 700 & 1.000 & 0.071 & \textbf{0.005} & 0.065 \\
Attention-only, parameter-matched & 600 & 1.000 & 0.074 & \textbf{0.004} & 0.066 \\
Recurrent channel ($\lambda{=}0$) & 600 & 1.000 & 0.070 & \textbf{0.004} & 0.064 \\
Recurrent channel ($\lambda{=}1$) & 600 & 1.000 & 0.067 & \textbf{0.004} & 0.062 \\
\bottomrule
\end{tabular}
\end{table}

\begin{table}[t]
\centering\small
\caption{Horizon and state fidelity by condition. $H_{\mathrm{ill}}(.05)$ is
the ply at which the illegal-move rate first crosses five percent,
interpolated between bucket midpoints because the bucketed read-off returns
the same bucket for every condition and so discriminates nothing.
Mean\,$\pm$\,s.d.\ over seeds where more than one was run.}

\begin{tabular}{lrcccc}
\toprule
Model & Params & seeds & occ.\ acc.\ @40--50 & illegal @40--50 & $H_{\mathrm{ill}}(.05)$ \\
\midrule
Attention-only & 7{,}345{,}920 & 3 & 0.555\,$\pm$\,0.003 & 0.176\,$\pm$\,0.002 & 26.1\,$\pm$\,0.0 \\
Attention-only, parameter-matched & 7{,}505{,}976 & 3 & 0.554\,$\pm$\,0.002 & 0.172\,$\pm$\,0.003 & 26.0\,$\pm$\,0.1 \\
Recurrent channel ($\lambda{=}0$) & 7{,}508{,}416 & 3 & 0.556\,$\pm$\,0.003 & 0.169\,$\pm$\,0.002 & 26.2\,$\pm$\,0.3 \\
Recurrent channel ($\lambda{=}1$) & 7{,}508{,}416 & 3 & 0.632\,$\pm$\,0.003 & 0.166\,$\pm$\,0.002 & 26.5\,$\pm$\,0.2 \\
\midrule
\quad read-out ablated at test & n/a & 3 & 0.559\,$\pm$\,0.005 & 0.149\,$\pm$\,0.005 & 27.2\,$\pm$\,0.4 \\
\midrule
\quad 6L/192 & 3{,}440{,}064 & 3 & 0.468\,$\pm$\,0.006 & 0.245\,$\pm$\,0.003 & 21.4\,$\pm$\,0.1 \\
\quad 8L/256 & 7{,}345{,}920 & 3 & 0.555\,$\pm$\,0.003 & 0.176\,$\pm$\,0.002 & 26.1\,$\pm$\,0.0 \\
\quad 12L/256 & 10{,}504{,}960 & 3 & 0.577\,$\pm$\,0.005 & 0.149\,$\pm$\,0.001 & 27.8\,$\pm$\,0.2 \\
\quad 8L/384 & 15{,}737{,}472 & 3 & 0.632\,$\pm$\,0.006 & 0.134\,$\pm$\,0.004 & 29.5\,$\pm$\,0.2 \\
\quad 12L/384 & 22{,}835{,}328 & 3 & 0.656\,$\pm$\,0.001 & 0.112\,$\pm$\,0.002 & 31.7\,$\pm$\,0.3 \\
\quad 12L/512 & 39{,}884{,}288 & 3 & 0.694\,$\pm$\,0.003 & 0.101\,$\pm$\,0.001 & 33.8\,$\pm$\,0.1 \\
\bottomrule
\end{tabular}
\end{table}

\begin{table}[t]
\centering\small
\caption{The scale ladder, three seeds per rung. $H_{\mathrm{ill}}$ is read
off behaviour, the ply at which the illegal-move rate crosses five percent.
$H_{\mathrm{fid}}$ is read off the representation, the depth at which the
probe first fails on the squares the chosen move touches more than half the
time, and so cannot be moved by hedging. Belief consistency is normalised by
each model's own legal-move recovery reference, because that reference rises with
size. Mean\,$\pm$\,s.d.\ over seeds.}

\begin{tabular}{lrccccc}
\toprule
Model & Params & seeds & $H_{\mathrm{ill}}(.05)$ & $H_{\mathrm{fid}}(.5)$ & belief\,/\,reference & mismatch \\
\midrule
6L/192 & 3{,}440{,}064 & 3 & 21.4\,$\pm$\,0.1 & 33.8\,$\pm$\,0.2 & 0.476\,$\pm$\,0.017 & 0.051\,$\pm$\,0.002 \\
8L/256 & 7{,}345{,}920 & 3 & 26.1\,$\pm$\,0.0 & 40.8\,$\pm$\,0.5 & 0.534\,$\pm$\,0.008 & 0.052\,$\pm$\,0.002 \\
12L/256 & 10{,}504{,}960 & 3 & 27.8\,$\pm$\,0.2 & 42.1\,$\pm$\,0.3 & 0.552\,$\pm$\,0.007 & 0.051\,$\pm$\,0.002 \\
8L/384 & 15{,}737{,}472 & 3 & 29.5\,$\pm$\,0.2 & 50.0\,$\pm$\,0.4 & 0.568\,$\pm$\,0.008 & 0.053\,$\pm$\,0.001 \\
12L/384 & 22{,}835{,}328 & 3 & 31.7\,$\pm$\,0.3 & 52.5\,$\pm$\,0.5 & 0.577\,$\pm$\,0.005 & 0.053\,$\pm$\,0.003 \\
12L/512 & 39{,}884{,}288 & 3 & 33.8\,$\pm$\,0.1 & 59.1\,$\pm$\,0.5 & 0.587\,$\pm$\,0.011 & 0.056\,$\pm$\,0.004 \\
\bottomrule
\end{tabular}
\end{table}

\begin{table}[t]
\centering\small
\caption{Planning error against how strictly we condition on state. Each
cell gives the median engine centipawn loss, the blunder rate, and the
sample size, over legal moves whose touched squares the probe recovers and
with at most the stated number of other squares wrong. Tightening the
conditioning lowers the loss, so part of what looks like degraded judgement
is state loss elsewhere on the board. Loss still grows with depth at fixed
conditioning, so part of it is not. Cells with fewer than thirty scored
positions are reported as n/a, and that scarcity is itself the finding:
at depth a nearly intact board is rare.}

\begin{tabular}{lcccccc}
\toprule
Ply & $\leq 0$ & $\leq 2$ & $\leq 4$ & $\leq 8$ & $\leq 16$ & any \\
\midrule
20--30 & n/a & 24 / 0.18 (115) & 24 / 0.21 (300) & 31 / 0.22 (300) & 44 / 0.31 (300) & 58 / 0.34 (300) \\
30--40 & n/a & n/a & 38 / 0.25 (85) & 62 / 0.36 (300) & 75 / 0.41 (300) & 91 / 0.47 (300) \\
40--50 & n/a & n/a & n/a & 67 / 0.41 (300) & 120 / 0.53 (300) & 111 / 0.52 (300) \\
\bottomrule
\end{tabular}
\end{table}

\begin{table}[t]
\centering\small
\caption{Domain 2, permutation tracking. Answer accuracy at the shallowest
and deepest depths against a chance rate of one in four, the probe's
recovery of the queried object at depth, and belief consistency beside its
mismatched control. The attention-only model solves the task well above
chance while its state stays at chance under the probe, so it is not
maintaining a linearly decodable assignment. The supervised channel is, and
its errors are coherent with it.}

\begin{tabular}{lccccc}
\toprule
Model & seeds & acc.\ shallow & acc.\ deep & probe deep & belief / control \\
\midrule
Attention-only & 3 & 1.000\,$\pm$\,0.000 & 0.589\,$\pm$\,0.010 & 0.263\,$\pm$\,0.010 & 0.256\,$\pm$\,0.005 / 0.249\,$\pm$\,0.003 \\
Recurrent channel ($\lambda{=}1$) & 3 & 1.000\,$\pm$\,0.000 & 0.989\,$\pm$\,0.001 & 0.983\,$\pm$\,0.001 & 0.701\,$\pm$\,0.017 / 0.254\,$\pm$\,0.009 \\
\bottomrule
\end{tabular}
\end{table}

\begin{figure}[t]
\centering
\includegraphics[width=6in]{decay.png}
\caption{Internal state and external behaviour decay together. (a) Occupied-square
state accuracy at the most decodable layer, against a majority predictor, a
label-permutation control, and a randomised-weight model of identical
architecture. (b) The illegal top-1 move rate and the probability mass the model
places on legal continuations.}

\end{figure}

Occupied-square decodability falls from 0.975 in the opening to
0.066 in the deepest bucket, while the illegal top-1 rate rises from
0.0004 to 0.319 and the probability mass on legal moves
falls from 0.993 to 0.455. At mid-game the trained model
is well clear of every trivial explanation, reaching 0.555 against
0.261 for the majority predictor and lower still for the
label-permutation and randomised-weight controls, which
Table~1 reports in full at every depth. There is a real state
representation, and it becomes progressively less recoverable by a linear
decoder. Section~
ef{sec:drift} tests whether that reflects loss or a moving
basis.

\subsection{Aggregate fidelity is the weaker signal}

The two curves fall together, and it would be easy to stop there and call them
one phenomenon. The relationship is real but weaker than that, and measuring it
carelessly overstates it.

An earlier version of this analysis compared a point-biserial correlation for the
aggregate predictor against a conditional probability for the move-touched
one. Those are different statistics, and the contrast flattered the conclusion: a
correlation against a heavily skewed binary outcome is attenuated, so the
aggregate measure looked uninformative when it was merely differently scaled. We
now score both the same way, as predictors of the same binary outcome on the same
samples, within depth so that neither can simply be tracking depth.

Aggregate board fidelity predicts an illegal next move with an AUC between
0.538 and 0.671. That is above chance and not
negligible. The same representation restricted to the two squares the move
touches reaches 0.592 to 0.687, and the two are
closest in the opening and furthest apart at depth: at ply 40 to 50 the aggregate
reaches 0.548 against 0.643. Fitted jointly
with depth as a covariate, the move-touched term takes a standardised weight
of 0.504 against 0.118 for the aggregate.

\begin{figure}[t]
\centering
\includegraphics[width=6in]{coupling.png}
\caption{(a) Within a fixed depth, the correlation between the number of
misremembered squares and playing an illegal move is small. (b) The same
representation, scored only on the squares the move actually uses, separates
illegal from legal moves sharply. Table~2 compares the two on equal
footing, as predictors of the same outcome.}

\end{figure}

The claim we can defend is therefore narrower than the one we first wrote, and it
is still a claim about a standard practice. Whole-board probe accuracy is a
usable but weak proxy for whether a world model is load-bearing, it degrades as a
proxy exactly where long-horizon behaviour matters most, and a measurement
restricted to the state an action consumes is several times more informative from
the same probe and the same activations. Stated as a conditional rather than a
ranking, the probe is wrong on the squares a move touches for
0.806 of illegal moves against 0.384 of legal ones.
The relevant question is not how much of the world the model has lost, but
whether it has lost the part it is about to use.

\subsection{Lost, or merely moved? Depth-local probes}

Every fidelity number so far comes from one decoder fitted across all depths, and
a falling number has been read as the state becoming less available. A single
global decoder cannot distinguish that from an alternative: the state might stay
equally present while its linear basis moves with depth, so that one fixed
decoder simply transfers badly. The reversible-cycle experiment holds the
position fixed but still uses the global probe, so it inherits the same
ambiguity.

We fit a separate linear probe inside each depth bucket and evaluate every probe
on every bucket.

\begin{table}[t]
\centering\small
\caption{Is the state lost, or is the representation rotating? Each row
is a linear probe fitted inside one depth bucket and evaluated on every
bucket; the pooled row is a probe fitted across all depths, as used
elsewhere in the paper. Depth-local probes do not hold up at depth either,
so the decline is not explained by a fixed decoder failing to transfer.
Local probes see roughly a sixth of the pooled training data, which makes
this a conservative test of the drift account.}

\begin{tabular}{lcccccc}
\toprule
Probe fitted at & ply 10 & ply 20 & ply 30 & ply 40 & ply 50 & ply 60 \\
\midrule
ply 10 & \textbf{0.910} & 0.776 & 0.521 & 0.307 & 0.203 & 0.151 \\
ply 20 & 0.835 & \textbf{0.789} & 0.695 & 0.540 & 0.348 & 0.201 \\
ply 30 & 0.640 & 0.736 & \textbf{0.671} & 0.610 & 0.546 & 0.412 \\
ply 40 & 0.253 & 0.591 & 0.636 & \textbf{0.553} & 0.486 & 0.416 \\
ply 50 & 0.106 & 0.341 & 0.527 & 0.514 & \textbf{0.430} & 0.346 \\
ply 60 & 0.097 & 0.241 & 0.415 & 0.497 & 0.447 & \textbf{0.340} \\
\midrule
pooled across depths & 0.907 & 0.804 & 0.694 & 0.591 & 0.506 & 0.409 \\
\bottomrule
\end{tabular}
\end{table}

The diagonal answers the question. A probe fitted for the deepest bucket and
tested there reaches 0.340, against 0.910 for the
shallowest bucket tested on itself. Depth-local probes collapse in almost the
same way the global probe does, 0.907 to 0.409, so
the decline is not an artifact of a fixed decoder failing to keep up. If the
representation merely rotated, a probe fitted at the relevant depth would recover
what a global probe misses, and it does not.

Basis change is nonetheless present in the off-diagonal. A probe fitted on the
shallowest positions reaches only 0.151 on the deepest ones, far
below what a probe fitted there achieves, so the encoding does move with depth.
It moves and it also thins, and the thinning is what the paper's curves track.

One caveat cuts in our favour rather than against. Each depth-local probe sees
roughly a sixth of the data the pooled probe sees, which is why the pooled probe
edges the diagonal at depth. More data would help the local probes, so this is a
conservative test of the drift account: it would have to overturn a gap that
already runs the wrong way for it.

\subsection{Which squares, not how many}

Move-touched fidelity scores two squares where the aggregate averages
sixty-four, and a statistic over two concentrated variables can look sharper than
one over sixty-four diluted ones for reasons unrelated to relevance. The control
that separates those explanations is a random pair.

\begin{table}[t]
\centering\small
\caption{Which squares carry the signal. Each row scores probe error
over a different set of squares as a predictor of an illegal top-1 move,
on identical samples. Intervals are bootstrapped over games rather than
positions, since positions within a game are not independent. Two random
occupied squares are a partial match to a move's endpoints, the origin being
occupied by the mover while the destination is usually empty or held by the
opponent, so this control equalises count and occupancy only in part. It
performs like the whole board; the squares the move actually touches do
not, so the advantage is about which squares rather than how many.}

\begin{tabular}{lc}
\toprule
Squares scored & AUC [95\% CI] \\
\midrule
whole board, 64 squares & 0.629 [0.623, 0.633] \\
two random squares & 0.539 [0.534, 0.542] \\
two random occupied squares & 0.631 [0.627, 0.636] \\
origin square only & 0.664 [0.660, 0.668] \\
destination square only & 0.637 [0.633, 0.642] \\
move-touched, origin and destination & \textbf{0.696} [0.693, 0.700] \\
\bottomrule
\end{tabular}
\end{table}

Two randomly chosen squares predict an illegal move at an AUC of
0.539, well below the whole board at 0.629. Two
squares drawn at random from those actually \emph{occupied} reach
0.631, statistically indistinguishable from the whole board.
The squares a move touches reach 0.696, with a game-clustered
interval that does not overlap either control. The advantage is therefore about
which squares are read, not how many, and it is not an artifact of averaging
fewer variables. That control equalises count and occupancy only in part: a
move's origin is occupied by the side to move, while its destination is usually
empty or held by the opponent, so a pair of occupied squares does not reproduce
the endpoint profile exactly. A fully matched control would draw one square
occupied by the mover and one matching the destination's true class.

Splitting the pair, the origin square carries more of the signal than the
destination, 0.664 against 0.637. That ordering is
what one would expect if the dominant failure is proposing a move from a square
the model has misremembered as occupied.

\subsection{The errors are coherent}

If state loss drives the failure, an invalid action should be valid in the state
the model believes it is in. We reconstruct that believed position from the probe
and re-score the model's own move in it, clearing castling rights so nothing is
licensed by an assumption we did not decode, and testing pseudo-legality because
a decoded board may lack a king.

\begin{figure}[t]
\centering
\includegraphics[width=6in]{belief.png}
\caption{The errors are coherent. (a) Against depth, restricted to buckets with
at least thirty illegal cases. (b) Pooled: an illegal move is far more likely to
be legal in the model's own believed board than under another position's belief
at equal depth, or than an arbitrary illegal move is under the model's own
belief. The dashed line is the recovery reference for this test, set by how often
genuinely legal moves are recovered as legal.}

\end{figure}

Of the illegal moves the model plays, 0.355 are legal in its own
believed position, with a game-clustered interval of [0.346,
0.364] over 13,802 illegal moves from 1,155 games.
Positions within a game are not independent, so the interval resamples whole
games. Two controls bound the alternatives. The same move scored
against a believed position decoded elsewhere at equal depth is legal
0.050 [0.046, 0.053] of the time, so the
effect is specific to this belief at this position rather than to plausible
boards generally. An arbitrary illegal move scored against the model's own
believed position is legal 0.007 [0.006,
0.009] of the time, so the believed board is not permissive. For reference, genuinely legal moves are recovered as legal in the believed board
0.654 of the time, which is below one because probe decoding is
lossy. We call this a legal-move recovery reference rather than a decoding
reference, and we do not divide by it to estimate an unobserved true rate: that
would assume probe sensitivity is the same on legal and illegal cases, which we
have not established. What the comparison licenses is that a substantial
identified subset of illegal moves is consistent with the decoded state, not that
most of them are.

We do not divide by that reference. Doing so would estimate an unobserved true
rate only under an assumption we have not tested, that the probe is equally
sensitive on legal and illegal cases. What the comparison licenses is narrower
and still substantial: a large identified subset of illegal moves is admissible
under the board this model appears to hold at this position, at
7.2 times the rate of the matched control. It identifies
state-consistent failure. It does not assign the remainder.

\subsection{Inducing a false belief changes the action}

Belief consistency is observational. It shows the model errs coherently with a
decayed state without showing that the state representation drives the action.
We therefore intervene. At a position where the side to move genuinely occupies a
square, we push the residual stream along the probe direction encoding "this
square is empty" and ask whether the model stops proposing moves from it.

Two safeguards decide whether such an experiment means anything, and both are
reported. The manipulation check asks whether the edit changed the decoded belief
at all: at the strength we use it does, in 1.000 of cases. The
norm-matched control applies a random direction of identical magnitude at the
same position, which separates belief from mere perturbation sensitivity.

Edit strength has to be calibrated rather than chosen, because a large enough
push degrades the forward pass and will move behaviour for reasons having nothing
to do with state. Table~11 reports the sweep. Below a quarter of
the residual-stream magnitude the edit does not reliably flip the belief; above
roughly one, the random control also starts destroying the output distribution,
so any effect there is damage rather than belief. We report the intervention at
0.50, inside the window where the manipulation succeeds and
the control does not.

\begin{figure}[t]
\centering
\includegraphics[width=6in]{patch_calibration.png}
\caption{Choosing the edit strength. Too small and the belief does not flip; too
large and a norm-matched random direction also destroys the output distribution,
so any effect there is damage rather than belief. We report inside the shaded
window.}

\end{figure}

At that strength, editing the state direction removes 0.934 of
the probability mass on moves from the affected square, decreasing it in
0.996 of 700 cases. The norm-matched random edit
removes 0.077 and decreases the mass in
0.484 of cases, which is indistinguishable from a coin flip.
The state representation is not a decorative correlate of the model's
computation. It is read.

\subsection{How much of the planning error is really memory}

The other half of the decomposition needs care, because the conditioning that
would isolate it cleanly is unattainable. Requiring the probe to recover the
exact board leaves no samples past the opening, for the reason already given.
Requiring only the squares the move touches leaves samples everywhere but does
not constrain the rest of the board, so state loss elsewhere is still counted as
a failure of judgement.

We therefore vary the strictness and read the trend (Table~7).
Tightening the conditioning lowers the loss substantially. At ply 30 to 40 the
median centipawn loss falls from 91 with the board
unconstrained to 38 when at most four other squares may be
wrong, and the blunder rate from 0.473 to
0.247. Roughly half of what a naive reading would call
degraded judgement is misattributed memory failure.

A limit of this analysis is worth naming, because it is not a nuisance but the
same fact the paper is about. The strictest cells are thinly populated at depth
however many games we scan, since a nearly intact board becomes rare exactly
where we would most like to condition on one. Cells with fewer than thirty scored
positions are reported as n/a rather than given a number.

The trend does not go to zero. Holding the conditioning fixed at four wrong
squares, the median loss still grows with depth, from
24 at ply 20 to 30 to 38 at ply 30
to 40. Judgement given intact state is not depth-invariant.

We take the honest reading. The identified state-consistent component is
substantial, stricter state conditioning reduces measured planning loss, and a
residual depth-dependent planning component remains. Calling memory dominant
would overstate what a lossy probe and small conditioned samples support. Our original framing of H4, a
clean separation in which an intervention moves one term and leaves the other
untouched, was too tidy for what the measurement supports. The decomposition
still earns its place, because it converts one uninterpretable curve into two
quantities with different remedies, but a practitioner should expect the terms to
be entangled rather than orthogonal, and should report the conditioning they used
when they quote either one.

\subsection{Scale}

The behavioural horizon is read off the illegal-move curve at the ply where it
first crosses five percent, interpolated between bucket midpoints because the
bucketed read-off returns the same bucket for every condition and so cannot
discriminate between them. A second horizon defined on the representation is
introduced in Section~7.16; both are tabulated together there.

Across the ladder the behavioural horizon grows with scale and grows slowly,
from 21.4 ply at 3,440,064 parameters to
33.8 at 39,884,288, three seeds per rung. Roughly an
order of magnitude in parameters buys about half again as much depth, which
is consistent with H2 and with the cost mismatch of Section~5:
scaling helps, and it does not buy back a missing mechanism cheaply.

The question that matters more for the reach of this paper is not whether the
horizon moves with scale, which it does, but whether the mechanism does.
Figure~6(b) and Table~6 report belief consistency at
every rung, three seeds each. The raw rate rises with size, but so does each
model's legal-move recovery reference, so the honest quantity is the ratio of the two.
Normalised that way the coherence of errors still rises monotonically, from
0.476 at 3,440,064 parameters to
0.587 at 39,884,288, with seed spreads of a few parts
in a thousand, while the mismatched control stays flat near
0.050 and the random-move control near 0.007. We show
the ratio as a descriptive normalisation, since the reference itself rises with
size, and not as an estimate of a hidden true proportion. What we have characterised strengthens with scale over
the range we can reach rather than washing out, which is evidence against its
being an artifact of the smallest models.

Both horizons order the models identically and both grow sublinearly in
parameters, the behavioural one from 21.4 to 33.8 ply
and the fidelity one from 33.8 to 59.1
(Table~6). Roughly an order of magnitude in parameters buys about
half again as much depth, which is consistent with H2 and with the cost mismatch
of Section~5: scaling helps, and it does not buy back a
missing mechanism cheaply.

\begin{figure}[t]
\centering
\includegraphics[width=6in]{scale.png}
\caption{(a) The horizon against model size. (b) Belief consistency and its
matched control against model size. The question is whether the mechanism is a
property of small models or survives scale.}

\end{figure}

\subsection{An intervention, evaluated by the diagnostic}

The diagnostic is only useful if it discriminates. We test a fixed-width gated
recurrent channel inserted after every $k$-th attention block, carrying
$z_t \in \mathbb{R}^{d_s}$ with $d_s \ll d$:
\[
u_t = W_{\mathrm{in}} h_t, \quad
\alpha_t = \sigma(W_g[h_t; z_{t-1}]), \quad
z_t = (1-\alpha_t)\odot A z_{t-1} + \alpha_t \odot \phi(u_t),
\]
read back as $h_t \leftarrow h_t + W_{\mathrm{out}} z_t$, with $A$ diagonal so
inference-time state is $O(1)$ per token, and an optional auxiliary loss
$\lambda\mathcal{L}_{\mathrm{state}}$ supervising $z_t$ to be the state. We
separate architecture ($\lambda=0$), supervision ($\lambda>0$), and read-out
(zeroing $W_{\mathrm{out}}$ at test time), against a parameter-matched
attention-only control whose feed-forward block absorbs the same parameter count.

The recurrent channel gives a clean and slightly uncomfortable answer.

With the auxiliary objective active, state decodability rises sharply, from
0.555 for the attention-only baseline to 0.632 at
matched parameters and matched tokens, a difference many times the seed spread.
The architecture without the objective, 0.556, is
indistinguishable from the baseline, as is the parameter-matched control at
0.554. So the gain comes from supervising the channel, not from
adding it, and not from the parameters it costs.

Behaviour barely follows. The illegal-move rate improves from
0.176 to 0.166, and the horizon from
26.1 to 26.5 ply. Set that beside the scale ladder:
the attention-only model at 15,737,472 parameters reaches almost
exactly the same state decodability, 0.632, and converts it into
a much larger behavioural gain, 0.134 and
29.5 ply.

Two models with the same measured world model, then, and very different
behaviour. This is the paper's thesis arriving from the other direction. In
Section~7.2 aggregate decodability predicted behaviour weakly
across positions within a model; here it fails to predict behaviour across
models. An auxiliary loss that makes the state linearly available does exactly
what it is asked to do, and being available is not the same as being used.

\paragraph{An ablation that improves behaviour.}
Zeroing the read-out at test time lowers the illegal-move rate, from
0.166 to 0.149, and lengthens the horizon, from
26.5 to 27.2 ply. Removing a component the model was
trained with should not help, so we measured what else it does. Language-model
loss rises when the read-out is removed, from 3.010 to
3.199 across seeds, and the read-out is not vestigial: it injects a
vector of 0.150 of the residual stream's own norm. The ablated model
is worse at predicting the continuation and simultaneously better at staying
legal, and the reason is visible in the distribution. It places more mass on
legal moves than the intact model, 0.647 against
0.615, while committing less confidently to any particular one.

The lesson is about our own observable rather than about the artifact. A model
can lower its illegal-move rate by hedging toward the legal-move prior, so the
illegal rate is not a monotone proxy for the quality of state maintenance, and
should be read alongside a likelihood. This does not disturb the coupling,
belief-consistency, or causal results, which are computed per position rather
than from a rate, but it does qualify the intervention comparison and the horizon
metric, and we flag it rather than leaving a reader to find it.

We therefore report H3 as supported only in its weakest form and H4 as not
supported in the form we predicted. We had expected an intervention that improved
state tracking to improve behaviour through it. What we observe is an
intervention that improves the measurement of state tracking far more than it
improves state tracking's consequences. Anyone using auxiliary state supervision
as a training signal, or probe accuracy as a progress metric, should treat that
gap as the default expectation rather than the surprise.

\begin{table}[t]
\centering\small
\caption{Calibrating the edit strength for the causal intervention. Strength is
in units of the residual stream's own magnitude. The usable window is where the
manipulation check succeeds while the norm-matched random control still behaves
like a coin flip. Values from \texttt{results/attn\_8L256\_s0\_patch\_sweep.json}.}

\begin{tabular}{cccccc}
\toprule
strength & belief flipped & state edit: mass after & state decreased
& random: mass after & random decreased \\
\midrule
0.10 & 0.652 & 0.046 & 0.980 & 0.076 & 0.540 \\
0.25 & 0.996 & 0.020 & 0.988 & 0.078 & 0.452 \\
0.50 & 1.000 & 0.005 & 0.996 & 0.065 & 0.484 \\
1.00 & 1.000 & 0.000 & 0.996 & 0.054 & 0.556 \\
2.00 & 1.000 & 0.000 & 0.996 & 0.025 & 0.712 \\
6.00 & 1.000 & 0.000 & 1.000 & 0.013 & 0.784 \\
\bottomrule
\end{tabular}
\end{table}

\subsection{A second domain, and what it took to build one}

We attempted a second oracle-bearing domain, a program updating integer variables
with interleaved queries, to test whether these findings are properties of
long-horizon state tracking or of chess. Four designs failed to leave chance, and
we report the sequence because the diagnosis is more useful than the intention.

The first placed a single query at the end of a sequence of random updates.
Training loss settled at the token-type entropy of the generator, which is the
value a model achieves by learning the format and nothing else, and the reason is
arithmetic: the only learnable signal was one answer token in roughly one
hundred and thirty, so state tracking carried under one percent of the gradient.
Interleaving queries raised that share and did not move accuracy. Reducing the
modulus and the number of variables did not move it either. A final version
separated a pure read-out query from a computed one, so that state tracking could
succeed without the arithmetic succeeding, and read-out accuracy still sat at
chance.

We draw two conclusions from the arithmetic attempts. The narrow one is that this task
family asks a small model to learn modular arithmetic before any state tracking
becomes observable, and that our compute did not reach it. The broader one is a
warning for anyone building synthetic state-tracking benchmarks: a task can be
perfectly well specified, carry an exact oracle, and still measure nothing,
because the quantity of interest contributes almost none of the training signal.
The loss curve diagnoses this cheaply. A model that has learned the format and no
content sits at the generator's token-type entropy, which is computable in
advance and worth computing before trusting such a benchmark. We report this in
place of an easy cross-domain result.

\paragraph{A domain that does work.}
Abandoning arithmetic entirely, we built the second domain around an assignment
instead. Four objects sit in four boxes. The sequence states the initial
assignment, then applies operations that permute it, and after every operation
asks where one object now is. A \textsc{swap} names two boxes and moves their
contents without naming the objects, so the affected object's location cannot be
retrieved by attending to its last mention and has to be maintained. A
\textsc{move} names the object, which gives an anchor and keeps the task
learnable. The state is the assignment of objects to boxes, exact at every
token, and the answer is a single symbol, so belief consistency here is exact
rather than estimated: a wrong answer either is or is not the box the model's own
probe says holds that object.

Two further design choices were forced on us, and both are worth recording. Our
first attempt used five objects. That makes the state an element of $S_5$, a
permutation-composition problem that is theoretically hard for the relevant
fixed-depth model classes under asymptotic expressivity assumptions
(Merrill et~al. 2024). That may have contributed, but a theorem about
asymptotic expressivity does not establish why one finite model failed to
optimise on one bounded-length dataset, and we do not claim group structure as
the cause. Four objects give $S_4$,
which is solvable. Our second attempt queried only occasionally, which left the
task all-or-nothing: with uniformly random updates and sparse queries a model
either tracks the permutation exactly or sits at chance, a structure that is hard
to learn from endpoint supervision. Querying after every operation supplies
partial credit, since a query one step after an operation is answerable directly,
and that is the foothold from which the model extends. Chess has this property
for free, because a plausible move needs only partial state.

With those two changes the task became measurable, and the outcome follows.

The outcome is not the clean replication we hoped for, and it is more
informative than one would have been (Table~8).

With the supervised recurrent channel the task is close to solved at every depth,
0.989 at the deepest bucket against a chance rate of
0.25, the probe recovers the queried object at
0.983, and belief consistency is 0.701
against a mismatched control of 0.254. The mechanism
reproduces.

The attention-only model behaves differently. It solves the task well above
chance, 0.589 at depth, while the probe recovers its state at
0.263, which is chance. Belief consistency is
0.256 against a control of 0.249, which is
no signal at all. The same probe procedure, on the same positions, recovers state
at 0.983 in the other model, so this is a fact about the
model and not a failure of the probe.

The reading we prefer is that this task does not force the model to maintain
state. Only one object is asked about at a time, so the answer can be resolved on
demand by tracing that object's history, and an explicit assignment carried
forward is unnecessary. Chess allows no such shortcut: predicting a legal move
requires the whole position, so the state must be maintained eagerly.

That gives a fourth lesson for anyone building these diagnostics, and it is the
one we would put first. A task can carry an exact oracle, be learnable, and still
fail to measure state maintenance, because it never obliges the model to
maintain anything. Belief consistency is only defined where a belief exists to
read, and whether one exists is a property of the task, not only of the
architecture.

\subsection{Is the measurement useful, or only diagnostic?}

Everything so far is a claim about mechanism. A reader is entitled to ask whether
the same measurement is worth anything operationally, and the only baseline that
matters is the model's own confidence. If an external state monitor merely
recovers what the model already knows about its own uncertainty, it is
interesting and redundant.

We test this as selective prediction. A system that wants to avoid acting on a
decayed state abstains on the positions it trusts least, and the question is what
to rank by. We compare the model's top-1 probability, the entropy of its
next-move distribution, and the probe's certainty about the squares the chosen
move touches. That last signal uses only the probe's own output, so it is
computable at inference from activations, the probe having been fitted against an
oracle offline.

One version of this experiment is not admissible and we report it only as a
bound. Ranking by whether the probe is \emph{wrong} on those squares requires the
oracle at inference, which is the very thing an abstention rule exists to avoid.
It is an upper bound on what a perfect state monitor could deliver, not a method.

The honest result is modest. Probe certainty alone is a weaker signal than the
model's own confidence, with an AUC of 0.593 against
0.628. Combined with confidence it is better than either alone: at
half coverage the illegal-move rate among retained positions falls to
0.124 against 0.141 for confidence and
0.131 for entropy, and in a joint fit both terms carry weight,
0.554 for confidence and 0.367 for probe certainty.
An external monitor therefore adds something the model does not self-report, and
it adds less than we hoped.

The gap to the bound is the interesting part. Ranking by oracle-verified
move-touched state reaches 0.093 at the same coverage, less than
half the base rate of 0.197. Most of the benefit an ideal state monitor
would deliver is not captured by the probe's own certainty, which says the
limitation is the quality of the state read-out rather than the premise. Better
read-outs, not a different signal, are the lever.

\subsection{Inference that respects the clustering}

Positions within a game are not independent: they share an opening, a player, and
an accumulating history. Intervals computed over positions therefore understate
uncertainty, and every interval in this section resamples whole games instead.
Section~7.5 already reports game-clustered intervals for belief
consistency on the full evaluation population. This section is a restricted
sensitivity analysis on a different population, 142,042 positions from
2,500 games of which 19,896 carry an illegal top-1 move, over
ply 10 to 80 where the depth buckets used by the predictor comparison are
defined. Estimates therefore differ slightly from the primary ones, and both are
reported rather than reconciled by choosing whichever is larger.

\begin{table}[t]
\centering\small
\caption{Inference clustered on the game. Intervals resample whole
games rather than positions, since positions within a game share an
opening and a history. The lower block compares nested models by held-out
log loss on a split by game, so each predictor's contribution is measured
out of sample. Lower is better.}

\begin{tabular}{lcc}
\toprule
Quantity & Estimate & 95\% CI (clustered) \\
\midrule
belief consistency & 0.3892 & [0.3822, 0.3976] \\
\quad mismatched-belief control & 0.0654 & [0.0620, 0.0687] \\
\quad random-move control & 0.0079 & [0.0068, 0.0093] \\
AUC, whole board & 0.6287 & [0.6233, 0.6333] \\
AUC, move-touched & 0.6961 & [0.6926, 0.6998] \\
\quad difference & +0.0674 & [+0.0618, +0.0732] \\
\midrule
Held-out log loss & \multicolumn{2}{c}{} \\
depth only & 0.3630 & \\
\quad $+$ whole board & 0.3595 & \\
\quad $+$ move-touched & 0.3451 & \\
\quad $+$ both & 0.3439 & \\
\quad $+$ both and confidence & 0.3340 & \\
\bottomrule
\end{tabular}
\end{table}

Belief consistency is 0.389 [0.382, 0.398], against
0.065 [0.062, 0.069] for a mismatched
belief and 0.008 [0.007, 0.009] for a
random illegal move. The intervals are far apart, and the separation does not
depend on treating positions as independent.

The two predictors of an illegal move separate as well. Whole-board error reaches
an AUC of 0.629 [0.623, 0.633] and
move-touched error 0.696 [0.693,
0.700]. Bootstrapping the difference as a paired quantity gives
+0.0674 [+0.0618, +0.0732], which excludes
zero.

A population-averaged logistic fit, with game as the cluster and an exchangeable
working correlation, gives coefficients with clustered standard errors on standardised
predictors: +0.586 (clustered s.e.\ 0.011) for move-touched
error against +0.161 (0.009) for the aggregate, with
+0.573 for low model confidence and +0.633 for depth. The
ordering survives clustered inference.

\paragraph{Out of sample, not only in sample.} Coefficients can be significant and
still add little. Splitting by game and fitting on seventy percent of games, the
held-out log loss falls from 0.3630 with depth alone to
0.3595 adding whole-board error, and to 0.3451
adding move-touched error instead. Using both gives 0.3439, and
adding the model's own confidence 0.3340. The
move-touched predictor buys several times what the aggregate buys out of sample,
and it buys most of it even when the aggregate is already present.

\paragraph{A held-out abstention rule.} Fitting and evaluating a selective
prediction rule on the same positions overstates it. Refitting on training games
and evaluating on held-out games, the illegal rate among retained positions at
half coverage is 0.054 for a rule combining the model's confidence
with probe certainty, against 0.070 for confidence alone and
0.083 for probe certainty alone, from a base rate of
0.141. An oracle-verified state monitor would reach
0.038, which remains the bound rather than a method: establishing it
at inference needs the ground truth the rule exists to avoid.

\subsection{A constructive edit: moving a piece in the model's belief}

Suppressing a square's occupancy removes the mass on moves from it, which is
consistent with a causal belief state and equally consistent with generic damage
to source-square features. A constructive edit separates the two. We take a
square $A$ the mover genuinely occupies and an empty square $B$, and edit toward
``$A$ is empty'' and ``$B$ holds the piece that was on $A$'' together. Generic
damage predicts that mass leaves $A$; only a belief the policy reads predicts
that mass arrives at $B$.

Both happen. Across 600 positions at the calibrated strength, mass on
moves out of $A$ changes by -0.0705 [-0.0810,
-0.0603] while a norm-matched random edit gives -0.0052
[-0.0092, -0.0006]. Mass on moves out of $B$, which
are illegal in the true position, changes by +0.0172 [+0.0134,
+0.0222] against +0.0010 [+0.0006,
+0.0016] for the random edit. Intervals are bootstrapped over
games.

Splitting the edit budget across two directions makes the manipulation less
reliable than the single-square version: $A$ decodes as empty in
0.512 of cases and $B$ as the transferred piece in 0.647.
Restricting to the 205 cases where both flips demonstrably
succeeded, the constructive effect strengthens, to +0.0291
[+0.0173, +0.0397], while the random control on $A$
becomes indistinguishable from zero. We read this as the policy consulting a
joined-up occupancy representation rather than a single feature, while noting it
establishes this for occupancy and not for the whole position.

\subsection{Off distribution: uniformly random legal play}

Every measurement so far uses held-out human games, which are in distribution.
Walker and Lyons (2026) report that a uniformly-random-legal-play split stays
discriminative where in-distribution performance saturates, so it is a fair test
of whether our findings lean on familiar opening theory and typical structure.

They do not. On 1,200 games of uniformly random legal play,
occupied-square accuracy falls from 0.889 to 0.112 and
the illegal rate rises from 0.013 to 0.350, both
steeper than on human games as expected for an unfamiliar distribution. Belief
consistency is 0.380 over 32,027 illegal moves, against a
random-move control of 0.008. That is slightly higher than the
in-distribution rate, not lower. Coherent error is not an artifact of the model
being on familiar ground.

\subsection{Transcript length against position difficulty}

The decay curves are read throughout as loss caused by composing more updates.
Depth is confounded with position difficulty: a position at ply 100 differs from
one at ply 10 in material, phase, and legal-move count. The curves alone do not
license the causal reading, and this test removes the confound.

We insert a reversible cycle, a knight out and back for each side, which returns
the position exactly while adding four plies of transcript. State is then read on
an identical position at several transcript lengths.

\begin{table}[t]
\centering\small
\caption{Transcript length with the position held fixed. A reversible
cycle of legal moves returns the position exactly while adding four plies
of history, so every row scores the same states at a different transcript
length. Intervals are bootstrapped over games.}

\begin{tabular}{lccc}
\toprule
Added plies & occupied-square acc. & 95\% CI & illegal rate \\
\midrule
0 & 0.736 & [0.722, 0.747] & 0.065 \\
4 & 0.715 & [0.702, 0.727] & 0.050 \\
8 & 0.706 & [0.692, 0.718] & 0.113 \\
12 & 0.696 & [0.683, 0.707] & 0.177 \\
\bottomrule
\end{tabular}
\end{table}

Fidelity falls monotonically, from 0.736 [0.722,
0.747] with no insertion to 0.696 [0.683,
0.707] after twelve added plies, on 400 positions whose
state is unchanged. The intervals at the two ends do not overlap. Transcript
length alone therefore causes measurable state loss.

The size matters as much as the sign. Twelve plies of pure length cost about four
points of accuracy, where the natural depth curve falls far further over its
range. Both factors contribute, and we do not attribute the whole curve to
composition depth. One caveat we cannot remove: the inserted moves are legal but
unlike training data, so part of the drop may be distribution shift rather than
length as such. The random-play result of Section~7.14 bounds
that concern without eliminating it.

\subsection{A policy-conditioned state-fidelity horizon}

Both horizons used so far are read off behaviour, and the read-out ablation
showed why that is fragile: a model can lower its illegal-move rate by hedging
toward the legal-move prior, without knowing the position any better. A horizon
defined on the representation is less exposed, because it does not read the
confidence of the output: a model that spreads its mass more evenly gains nothing
from doing so. It is not policy-independent, and we do not claim it is. The
squares it scores are those the selected move touches, and that move comes from
the output distribution, so a policy change can alter which squares are scored
even with the state representation unchanged. It is a policy-conditioned
state-fidelity horizon. A genuinely policy-independent comparison would score a
fixed set of components, for instance the recorded continuation move, every
oracle-legal move averaged uniformly, or a candidate set shared across models.

We define $H_{\mathrm{fid}}$ as the depth at which the probability that the probe
is wrong on the squares the model's chosen move touches first crosses one half.
It uses the same probe and the same activations as everything else, and it is
immune to hedging: a model that spreads mass over legal moves gains nothing,
since the measure asks what it knows, not what it emits.

Read this way the ladder is orderly. The smallest model reaches
33.8 ply, the baseline 40.8, and the largest
59.1. The behavioural horizon orders the same models the same way,
21.4, 26.1 and 33.8 ply, on a shorter
scale because it asks a harder question of the model. Move-touched error at
the shallowest depth runs from 0.207 in the smallest model to
0.078 in the largest. We report both horizons throughout, and
prefer this one where the comparison is between architectures rather than between
depths, since that is where hedging is most likely to confound a behavioural
reading.

\subsection{A degenerate metric, declared}
Our horizon was pre-registered on exact-position fidelity, the probability that
all sixty-four squares are simultaneously recovered. That criterion is already
below its threshold in the very first bucket, so the pre-registered horizon
evaluates to zero ply for every condition and cannot discriminate between them.
The underlying quantity falls to 0.379 in the opening and to
0.000 by ply 20 to 30, after which exact states essentially do not
occur, which is also why conditioning on them yields no samples at depth. We therefore report horizons on
occupied-square accuracy and on the illegal-move rate. The revision was made
after inspecting baseline data and before evaluating any intervention, and every
definition is applied unchanged across conditions. We flag it rather than
presenting the revised metric as the original plan, and we note that anyone
measuring state tracking with an all-or-nothing criterion will obtain an
unusable number for the same reason.

\subsection{What the predictions came to}

Pre-registration is only worth the words if the verdicts are stated, including
the ones that went against us.

\textbf{H1, supported.}  State fidelity declines with depth, from
0.975 to 0.066 occupied-square accuracy, clear of a
majority predictor, a label-permutation control, and a randomised-weight model
(Table~1, Figure~2). The behavioural horizon,
26.1 ply, sits well inside the context window.

\textbf{H2, supported.}  The horizon grows slowly with scale across six rungs at
three seeds each, from 21.4 ply at 3,440,064 parameters
to 33.8 at 39,884,288 (Section~7.8).

\textbf{H3, supported only weakly.}  The supervised recurrent channel raises state
decodability substantially and the horizon by 26.5 against
26.1 ply. The architecture without supervision changes neither.

\textbf{H4, not supported as predicted.}  We expected an intervention that improved
state tracking to improve behaviour through it. Instead the intervention improved
the measurement far more than its consequences, matching a model twice its size
on decodability while capturing a fraction of the behavioural gain
(Section~7.9). We designed H4 to be able to fail and it failed, in a
direction that supports the paper's main claim rather than undermining it.

\textbf{H5, partially addressed.}  The natural-language transfer test our protocol
specifies was not run and is absent rather than approximated
(Section~9). A second oracle-bearing domain, reached only after
several failed designs, tests the weaker claim
(Section~7.10).

Three results were not predicted at all. The layer profile ran opposite to our
expectation (Section~5). Ablating the recurrent read-out
improved behaviour while worsening prediction, which turned out to say something
about our own observable rather than about the artifact
(Section~7.9). And the standard whole-board measure of a world model
proved a weaker predictor of behaviour than the same measure restricted to the
state an action consumes, by a margin that grows with depth
(Section~7.2).

If we could keep only one result it would be belief consistency together with the
causal edit, which say that these models act coherently on a state they are
losing. The measurement-placement result is a practical corollary: it tells you
where to look, and it matters most exactly where long-horizon behaviour is
hardest.

\section{Discussion}

The practical content of this paper is a change in what to measure. Reporting
that a probe recovers a world model at some aggregate accuracy is weak evidence
that the model is load-bearing: aggregate accuracy predicts the validity of the
next action, and it predicts it poorly, and it gets worse with depth. Measuring
the same representation on the components an action depends on is substantially
more informative from the same activations. Any team instrumenting a system
for long-horizon reliability should condition on the state the action uses.

The belief-consistency result reframes the failure. A model that plays a coherent
move in a misremembered position is not behaving erratically, and treatments that
assume erratic behaviour are aimed at the wrong target. Sampling several chains
and voting may offer limited benefit when independently sampled continuations
share the same state distortion, though we did not test repeated sampling or how
far drift is correlated across samples. A
verifier that scores outputs is well placed to catch the error and badly placed
to prevent it, and it will pay that cost on every query for as long as the system
runs.

The protocol generalises past chess. Anywhere an external checker can reconstruct
the intended intermediate state, program execution traces, spreadsheet
recalculation, structured database updates, ledger reconciliation, the same
measurements are available. The requirement is not a game. It is an oracle.

We also treat the negative results as part of the contribution. A predicted layer
profile came out backwards. A pre-registered metric turned out to be degenerate.
A first design of a second domain turned out to be unlearnable, for a reason that
was diagnosable from the loss curve. Each changed a decision, and each is
reported where it did so.

\paragraph{What this implies for practice.}
The decomposition tells a team which remedy fits the failure they have. Where
state-tracking error dominates, an inference-time verifier is an expensive way to
buy back something the architecture gave away, and it is a recurring per-query
cost rather than a one-off. Where planning error dominates, a verifier is exactly
right. Reporting only the sum leaves a practitioner unable to tell which world
they are in, and the two worlds recommend opposite purchases. The measurement
needed to tell them apart is a rules engine and a linear probe, not privileged
access to a model at scale, so it can be run in-house on a system a team already
depends on. For the reader of a long generated answer, who is rarely in a
position to audit step twenty-eight, the practical value is a stated depth beyond
which the output should not be trusted without checking.

\section{Limitations}

\textbf{L1}  \textbf{Scale.} Our models are small and trained on a modest corpus,
far below the chess-specific literature and any frontier system. Undertrained
models are easier to study because they fail within a measurable range, which is
what makes the instrument work, but the horizon values describe these models.
Section~7.8 estimates the trend rather than assuming it away.

\textbf{L2}  \textbf{One domain family.} Chess carries an exact oracle by
construction, which is what makes it measurable and also what makes it
unrepresentative of open-ended reasoning. The second domain we built does reach a
measurable regime (Section~7.10), but only its supervised
condition maintains a decodable state, so the attention-only replication we
wanted is not available there. The central findings rest on one instrument and
should be read that way.

\textbf{L3}  \textbf{Probe decoding is lossy.} The believed board is an
approximation, which is why we report the legal-move recovery rate as an explicit
recovery reference rather than treating the decoded state as ground truth.

\textbf{L4}  \textbf{Pseudo-legality.} Belief consistency is tested with
pseudo-legal moves, since a decoded board may lack a king and make check tests
ill-defined. This is a permissive criterion, which the random-move control is
designed to bound.

\textbf{L5}  \textbf{Scope of the causal claim.} The intervention establishes that
the decoded state direction is read by the policy at the layer tested. It does
not establish which downstream computation consumes it, and a single-direction
edit is a coarse instrument.

\textbf{L6}  \textbf{A metric was revised after seeing data.} Declared in
Section~7, applied uniformly, reported alongside the original.

\textbf{L7}  \textbf{No natural-language transfer test.} Establishing that these
findings hold for mathematical or textual reasoning needs matched general-corpus
pre-training, which our compute did not cover. It is absent rather than
approximated.

\section{Avenues of Future Research}

\textbf{FR1}  Extend the scale ladder far enough to estimate how horizon grows with
parameters and data, and test whether the belief-consistency rate is
scale-invariant even where the horizon moves.

\textbf{FR2}  Replicate the decomposition in domains that were not designed as
benchmarks, such as program execution or ledger reconciliation, where an oracle
exists but the data is not constructed.

\textbf{FR3}  Run the natural-language transfer test with matched pre-training, the
experiment that decides whether this line of work reaches practice.

\textbf{FR4}  Move from observation to intervention, editing move-touched state
directions and measuring whether the induced action follows the induced belief.

\textbf{FR5}  Develop state supervision that does not require an oracle, for example
from execution traces or agreement between sampled continuations.

\textbf{FR6}  Use move-touched fidelity as a live confidence signal, abstaining or
escalating when the state an action depends on is poorly determined.

\textbf{FR7}  Pre-register horizon metrics on pilot data, so metric selection is
settled before confirmatory runs rather than declared afterwards.

\section{Conclusion}

Long-horizon behavioural failure can arise from degradation of the task state,
from poor action selection given the state that is available, or from failures
that remain unresolved under a lossy decoder. These are not cleanly separable
with the instruments we have, and we have not separated them.

What we can report is narrower and, we think, still worth having. In search-free
chess models, whole-board probe fidelity is a weaker predictor of an invalid next
action than fidelity on the endpoints of the move actually selected, and the gap
widens with depth. A substantial identified subset of invalid actions is
admissible under the board that model appears to hold at that position, at many
times the rate of a mismatched-belief control and of an arbitrary invalid move.
Targeted suppression and transfer edits show that the corresponding occupancy
representation influences the policy rather than merely accompanying it. The
identified share grows across the range of model sizes we can train, and survives
a shift to uniformly random legal play.

These measurements identify state-consistent failure. They do not partition all
failures, and the errors they leave unidentified are unresolved rather than
attributed to reasoning. The practical consequence is a change in what to
measure: a world model reported as an aggregate accuracy tells a practitioner
less than the same representation scored on the state the next action will
consume, and the difference matters most exactly where long-horizon behaviour is
hardest to check.

\section*{Declarations}

\paragraph{Disclosure statement.}

The authors report no conflict of interest.

\paragraph{Funding.}

This research received no specific grant from any funding agency in the public,
commercial, or not-for-profit sectors.

\paragraph{Data availability statement.}

The training and evaluation code, the measurement pipeline, and the result files
from which every reported value is generated are available at
\url{https://github.com/MohitPratapSinghRathore/content-rot}.

\paragraph{AI usage declaration.}
During the preparation of this manuscript the authors used a generative AI tool
to assist with language editing and refinement of the text. The authors reviewed
and edited all output and take full responsibility for the content of the paper.
Generative AI was not used to generate or analyse the study's data, which were
produced entirely by the evaluation pipeline described in the paper.

\appendix

\section{Experimental detail}

\paragraph{Data.} Public Lichess archives, standard chess only, both players rated
at least 1600, games between 30 and 160 ply, deduplicated by a hash of the move
sequence. Moves are tokenised one per move in UCI form, giving a vocabulary of
1,926 including padding and a start symbol; there is no sub-move
tokenisation and no board rendering in the input. Sequences are padded to 161
tokens and truncated at 160 ply. Splits are by game: 96,000 for
language-model training, 12,000 for probe fitting,
12,000 held out for evaluation.

\paragraph{Oracle.} Every game is replayed with a rules engine, recording after
each token the occupancy of all 64 squares as one of 13 classes, the side to
move, and castling rights. Replay completeness is verified before any model is
measured.

\paragraph{Models.} Pre-norm decoder-only Transformers with learned positional
embeddings, GELU feed-forward of width four times the model dimension, and no
dropout. The grid spans 6 to 12 layers and widths 192 to 512, with 6 or 8 heads.
Training uses AdamW, $\beta = (0.9, 0.95)$, weight decay 0.01, a one-cycle
schedule peaking at $3\times10^{-4}$ with 5 percent warmup, batch size 32, 12000
steps, and mixed precision with gradient clipping at 1.0. The final checkpoint is
used; no checkpoint selection is performed. Seeds 0, 1 and 2 control
initialisation, batch order, and any auxiliary sampling.

\paragraph{Probes.} One linear head per layer maps the residual stream to 64
squares by 13 classes, trained with cross-entropy on the probe split only, AdamW
at $10^{-3}$, two epochs, batch size 48. The reported layer is the one with the
highest occupied-square accuracy, selected on the probe split rather than on the
evaluation split. Controls are a per-square per-depth majority predictor, a
label-permutation control pairing activations with another game's state at equal
depth, and a randomised-weight model of identical architecture whose probes are
fitted the same way.

\paragraph{Believed boards.} A decoded board is built by placing the argmax class
at each square, setting side to move from ply parity, and clearing castling
rights so that castling cannot be licensed by an assumption we did not decode.
Legality is tested as pseudo-legality, since a decoded board may lack a king and
make check tests ill-defined. Moves that fail to parse are counted as illegal.

\paragraph{Engine scoring.} Stockfish 17 at fixed depth 10. Loss is the drop from
the best available score to the score after the played move, both from the
mover's point of view, with mate scores mapped to 10000. The blunder threshold is
100 centipawns, fixed in advance.

\paragraph{Interventions.} Edits are added to the residual stream at the output of
the probed block, at the final position only. The direction is the difference of
probe decoder rows for the target and current classes, normalised and scaled to
$\alpha$ times the residual stream's own root-mean-square magnitude times the
square root of the model dimension. The reported $\alpha$ is 0.50,
chosen from the sweep in Table~11 as the smallest value that
reliably flips the decoded belief while leaving a norm-matched random edit at
chance. Random controls are drawn isotropically and rescaled to the same norm.

\paragraph{Recurrent channel.} A gated diagonal recurrence of width 64 inserted
after every fourth block, with the read-out matrix zero-initialised so the module
begins as the identity. The auxiliary objective, when active, is cross-entropy
from the channel to the oracle state with weight 1.0. The parameter-matched
control widens the feed-forward multiplier to 4.156 to absorb the same parameter
count.

\paragraph{Second domain.} Four objects in four boxes, an explicit initial
assignment, then 24 operations drawn evenly between a swap of two named boxes and
a move of a named object to a named box, with a query after every operation
asking where one object is. Answers are weighted five times in the loss. Details
of the six designs that failed before this one are in Section~7.10.

\section{Deviations from the pre-registration}

We record every deviation rather than presenting the final protocol as the
original one.

\begin{enumerate}[nosep]
\item The horizon was pre-registered on exact-position fidelity, which proved
degenerate. Horizons on occupied-square accuracy and on the illegal-move rate
were added after inspecting baseline data and before evaluating any intervention,
and applied unchanged to all conditions.
\item The bucketed horizon was replaced by an interpolated crossing point,
because the bucketed read-off returned the same bucket for every condition.
\item Planning error was pre-registered as conditional on exact state recovery,
which yields no samples past the opening. It is reported conditional on
move-touched recovery, with a strictness ladder showing what the choice costs.
\item The aggregate and move-touched comparison was first made with
incommensurable statistics, a correlation against a conditional probability. Both
are now scored as predictors of the same outcome.
\item The layer-profile prediction was falsified and is retained as such.
\item The second domain required seven designs. All are reported.
\end{enumerate}

\section{Research questions, design principles, and implications}

This work began as a design-science exercise, and the framing is recorded here
rather than in the body, where an empirical narrative reads better without it.

\paragraph{Research questions and objectives.}
\textbf{RQ1}, can long-horizon accuracy loss be separated into state tracking and
planning, is addressed by \textbf{RO1}, defining both and giving an estimator for
each (Sections~4 and 6). \textbf{RQ2}, does internal
state loss explain behavioural failure or merely accompany it, is addressed by
\textbf{RO2}, testing the link per token at the level of the state an action
consumes (Sections~7.2 and 7.5). \textbf{RQ3}, how far
does a search-free model reason before its state stops being recoverable and how
does that respond to scale, is addressed by \textbf{RO3}, measuring horizons
across a seeded ladder against trivial baselines (Section~7.8).
\textbf{RQ4}, does the diagnostic discriminate between architectural
interventions, is addressed by \textbf{RO4}, evaluating a recurrent state channel
with parameter-matched and ablated controls (Section~7.9).

\paragraph{Design principles.}
\textbf{DP1}, ground the diagnostic in an external oracle rather than model
self-report, realised by rules-engine labels at every token. \textbf{DP2},
separate the failure modes before attributing any effect, realised by scoring
planning only where the move is legal and the relevant state is independently
verified. \textbf{DP3}, establish what a trivial predictor achieves first,
realised by majority, label-permutation and randomised-weight controls on every
decodability number. \textbf{DP4}, measure the state the action depends on rather
than the state in aggregate, realised by move-touched fidelity and its
random-pair controls. \textbf{DP5}, test coherence by intervention and not only
by correlation, realised by the calibrated suppression and transfer edits.

\paragraph{Implications for individuals, organisations, and society.}
For the reader of a long generated answer, who is rarely placed to audit step
twenty-eight, the practical value is a stated depth beyond which output should
not be trusted unchecked, and the uncomfortable finding that failure arrives
fluent rather than obviously broken. For organisations, the decomposition
indicates which remedy fits: a recurring per-query verifier cost suits a policy
problem, a one-off structural change suits a memory problem, and reporting only
the sum leaves a team unable to tell which they have. Strategically, the
measurement needs a rules engine and a linear probe rather than privileged access
at scale, so it can be owned and audited in-house rather than taken on trust.
For society, diagnosing a failure is cheaper than paying inference-time compute
to mask it, and remedies that cost a small fraction of parameters travel to
on-device and public-sector deployments that cannot buy deliberation per request.

\section*{References}

Balogh, A. and Jelasity, M. (2026). Verification of the implicit world model in a generative model via adversarial sequences. In {\em International Conference on Learning Representations (ICLR)}.

Gu, A. and Dao, T. (2024). Mamba: Linear-time sequence modeling with selective state spaces. In {\em First Conference on Language Modeling (COLM)}.

Gu, A., Goel, K., and R{\'e}, C. (2022). Efficiently modeling long sequences with structured state spaces. In {\em The Tenth International Conference on Learning Representations (ICLR)}.

Harang, R., Naradowsky, J., Gujju, S.~V., and Miyao, Y. (2025). Tracking world states with language models: State-based evaluation using chess. In {\em ICML Workshop on Assessing World Models}. Spotlight; arXiv:2508.19851.

Hewitt, J. and Liang, P. (2019). Designing and interpreting probes with control tasks. In {\em Proceedings of the 2019 Conference on Empirical Methods in Natural Language Processing and the 9th International Joint Conference on Natural Language Processing (EMNLP-IJCNLP)}, pages 2733--2743. Association for Computational Linguistics.

Hsieh, C.-P., Sun, S., Kriman, S., Acharya, S., Rekesh, D., Jia, F., Zhang, Y., and Ginsburg, B. (2024). {RULER}: What's the real context size of your long-context language models? In {\em First Conference on Language Modeling (COLM)}.

Karvonen, A. (2024). Emergent world models and latent variable estimation in chess-playing language models. In {\em First Conference on Language Modeling (COLM)}.

Karvonen, A., Wright, B., Rager, C., Angell, R., Brinkmann, J., Smith, L., Verdun, C.~M., Bau, D., and Marks, S. (2024). Measuring progress in dictionary learning for language model interpretability with board game models. In {\em Advances in Neural Information Processing Systems 37 (NeurIPS 2024)}.

Li, K., Hopkins, A.~K., Bau, D., Vi{\'e}gas, F., Pfister, H., and Wattenberg, M. (2023). Emergent world representations: Exploring a sequence model trained on a synthetic task. In {\em The Eleventh International Conference on Learning Representations (ICLR)}.

Li, Q. and Jiang, W. (2026). Tracking vs. deciding: The dual-capability bottleneck in searchless chess transformers. arXiv:2603.29761.

Lightman, H., Kosaraju, V., Burda, Y., Edwards, H., Baker, B., Lee, T., Leike, J., Schulman, J., Sutskever, I., and Cobbe, K. (2024). Let's verify step by step. In {\em The Twelfth International Conference on Learning Representations (ICLR)}.

Liu, N.~F., Lin, K., Hewitt, J., Paranjape, A., Bevilacqua, M., Petroni, F., and Liang, P. (2024). Lost in the middle: How language models use long contexts. {\em Transactions of the Association for Computational Linguistics}, 12:157--173.

Meng, K., Bau, D., Andonian, A., and Belinkov, Y. (2022). Locating and editing factual associations in {GPT}. In {\em Advances in Neural Information Processing Systems 35 (NeurIPS 2022)}.

Merrill, W. and Sabharwal, A. (2024). The expressive power of transformers with chain of thought. In {\em The Twelfth International Conference on Learning Representations (ICLR)}.

Merrill, W., Petty, J., and Sabharwal, A. (2024). The illusion of state in state-space models. In {\em Proceedings of the 41st International Conference on Machine Learning (ICML)}, volume 235, pages 35492--35506.

Nanda, N., Lee, A., and Wattenberg, M. (2023). Emergent linear representations in world models of self-supervised sequence models. In {\em Proceedings of the 6th BlackboxNLP Workshop: Analyzing and Interpreting Neural Networks for NLP}, pages 16--30. Association for Computational Linguistics.

Pereira, D. and Zuidema, W. (2026). Transformers struggle to use their emergent world models: Revisiting the tower of hanoi, and the illusion of thinking. arXiv:2608.07077.

Ruoss, A., Del{\'e}tang, G., Medapati, S., Grau-Moya, J., Wenliang, L.~K., Catt, E., Reid, J., Lewis, C.~A., Veness, J., and Genewein, T. (2024). Amortized planning with large-scale transformers: A case study on chess. In {\em Advances in Neural Information Processing Systems 37 (NeurIPS 2024)}.

Toshniwal, S., Wiseman, S., Livescu, K., and Gimpel, K. (2022). Chess as a testbed for language model state tracking. In {\em Proceedings of the AAAI Conference on Artificial Intelligence}, volume~36, pages 11385--11393.

Vaswani, A., Shazeer, N., Parmar, N., Uszkoreit, J., Jones, L., Gomez, A.~N., Kaiser, {\L}., and Polosukhin, I. (2017). Attention is all you need. In {\em Advances in Neural Information Processing Systems 30 (NIPS 2017)}.

Walker, B. and Lyons, T. (2026). Chess-world-model: A 10m-game benchmark for exact state tracking from chess move sequences. arXiv:2605.30100.

Wang, X., Wei, J., Schuurmans, D., Le, Q., Chi, E.~H., Narang, S., Chowdhery, A., and Zhou, D. (2023). Self-consistency improves chain of thought reasoning in language models. In {\em The Eleventh International Conference on Learning Representations (ICLR)}.

Wei, J., Wang, X., Schuurmans, D., Bosma, M., Ichter, B., Xia, F., Chi, E.~H., Le, Q.~V., and Zhou, D. (2022). Chain-of-thought prompting elicits reasoning in large language models. In {\em Advances in Neural Information Processing Systems 35 (NeurIPS 2022)}.

Yao, S., Yu, D., Zhao, J., Shafran, I., Griffiths, T.~L., Cao, Y., and Narasimhan, K. (2023). Tree of thoughts: Deliberate problem solving with large language models. In {\em Advances in Neural Information Processing Systems 36 (NeurIPS 2023)}.

\end{document}
